\NumberedItem{5}
\HistoricTitle{\textsc{Gauss} an \textsc{Laplace}.}
\begin{flushright}\bfseries Göttingen ce 30. Janvier 1812\end{flushright}
\noindent\hspace*{.22\linewidth}{\bfseries Monsieur}\par\smallskip
\noindent\hspace*{\parindent}Je vous dis mille remercimens pour les deux memoires que vous m’avés fait l’honneur de m’envoyer[\textsuperscript{***})] et que j’ai reçus dans ces jours. Les fonctions que vous y traités aussi bien que les questions de probabilités, sur lesquelles vous preparés un grand ouvrage ont un grand attrait pour moi, quoique j’aie peu travaillé moimeme sur les dernieres. Je me rappelle pourtant d’un probleme curieux, duquel je me suis occupé il y a 12 ans, mais lequel je n’ai pas reussi alors à resoudre à ma satisfaction[$\dagger$)]. Peut être daignerés vous
\par\medskip\noindent\rule{.24\linewidth}{.4pt}\par\smallskip
{\small
\noindent\hspace*{\parindent}[\textsuperscript{***}) Es handelt sich wohl um die beiden folgenden Abhandlungen von \textsc{Laplace}: \textit{Mémoire sur les approximations des formules qui sont fonctions de très grands nombres et sur leur application aux probabilités}, Mémoires de l’Académie des Sciences \textbf{10} (1809), Paris 1810, S. 353 und \textit{Mémoire sur les intégrales définies et sur leur application aux probabilités etc.}, ebenda \textbf{11} (1810), Paris 1811, S. 279.]
\par\noindent\hspace*{\parindent}[$\dagger$) Siehe die Tagebuchaufzeichnung Nr. 113 vom 25. Oktober 1800.]
\par}
\par\noindent\hfill 47\textsuperscript{*}\par
\clearpage\SourceLeaf{379}\ArithmeticHead{372}{L}{ANALYSIS.\quad BRIEFWECHSEL.}
\noindent vous en occuper quelques momens: dans ce cas je suis sur que vous trouverés une solution plus complète. Le voici. Soit $M$ une quantité inconnue entre les limites 0 et 1, pour laquelle tou[te]s les valeurs sont ou egalement probables ou plus ou moins selon une loi donnée: qu’on la suppose convertie en une fraction continue
\[
M=\frac{1}{a'+\dfrac{1}{a''+\text{etc.}}}
\]

\indent Quelle est la probabilité, qu’en s’arretant dans le developpement à un terme fini, $a^{(n)}$, la fraction suivante
\[
\frac{1}{a^{(n+1)}+\dfrac{1}{a^{(n+2)}+\text{etc.}}}
\]
\noindent soit entre les limites 0 et $x$? Je la designe par $P(n,x)$ et j’ai en supposant pour $M$ toutes les valeurs également probables
\[
P(0,x)=x;
\]
\noindent $P(1,x)$ est une fonction transcendente dependante de la fonction
\[
1+\frac{1}{2}+\frac{1}{3}+{.}\,{.}\,{.}+\frac{1}{x}
\]
\noindent que \textsc{Euler} nome inexplicable et sur la quelle je viens de donner plusieurs recherches dans un memoire presenté à notre société des sc[iences] qui sera bientot imprimé[\textsuperscript{*})]. Mais pour les cas ou $n$ est plus grand, la valeur exacte de $P(n,x)$ semble intraitable. Cependant j’ai trouvé par des raisonnemens tres simples que pour $n$ infini on a
\[
P(n,x)=\frac{\log(1+x)}{\log 2}.
\]
\noindent Mais les efforts que j’ai fait lors de mes recherches pour assigner
\[
P(n,x)-\frac{\log(1+x)}{\log 2}
\]
\noindent pour une valeur tres grande de $n$, mais pas infinie, ont été infructueux.
\par\noindent\makebox[\linewidth][s]{{.}\hfill{.}\hfill{.}\hfill{.}\hfill{.}\hfill{.}\hfill{.}\hfill{.}\hfill{.}\hfill{.}\hfill{.}\hfill{.}\hfill{.}\hfill{.}\hfill{.}\hfill{.}\hfill{.}\hfill{.}\hfill{.}\hfill{.}\hfill{.}\hfill{.}\hfill{.}\hfill{.}\hfill{.}\hfill{.}\hfill{.}}\par
\medskip\noindent\rule{.24\linewidth}{.4pt}\par\smallskip
{\small
\noindent\hspace*{\parindent}[\textsuperscript{*}) \textit{Disquisitiones circa seriem etc.}, der Gesellschaft der Wissenschaften vorgelegt am 30. Januar 1812, Werke III, S. 123; siehe insbesondere S. 154 ff.]
\par}
\clearpage\SourceLeaf{380}\ArithmeticHead{373}{R}{GAUSS AN LAPLACE.}
\noindent\hspace*{\parindent}J’ai fait usage de la methode des moindre[s] carrés depuis l’an 1795 et je trouve dans mes papiers, que le mois de Juin 1798 est l’epoque où je l’ai rapprochée aux principes du calcul des probabilités: une note la dessus se trouve dans un journal que j’ai tenu sur mes occupations mathematiques depuis l’an 1796, et que j’ai montré dans ces jours à Mr. \textsc{de Lindenau}[\textsuperscript{*})].

Cependant mes applications \textls[120]{frequentes} de cette methode ne datent que des l’année 1802, depuis ce tems j’en fait usage pour ainsi dire tous les jours dans mes calculs astronomique[s] sur les nouvelles planetes. Comme je m’etois proposé depuis ce tems de reunir toutes les methodes dont je me suis servi dans un ouvrage etendu, (que j’ai commencé 1805 et dont le Manuscrit d’abord en allemand étoit achevé en 1806, mais lequel à la priere de Mr. \textsc{Perthes} j’ai traduit depuis en latin; l’impression a commencé en 1807 et n’est finie qu’en 1809[\textsuperscript{**})]), je ne me suis pas haté d’en publier un morceau detaché, ainsi Mr. \textsc{Legendre} m’est prevenu. Au reste j’avais deja communiqué cette meme methode, beaucoup avant la publication de l’ouvrage de M. \textsc{Legendre}[\textsuperscript{***})], à plusieurs personnes, entre autres à Mr. \textsc{Olbers} en 1803 qui certainement doit se le rappeler. Ainsi, pouvais je dans ma \textit{theorie des mouvement\textup{[s]} des planetes}, parler de la methode des moindre[s] quarrés, dont j’avais fait depuis 7 ans mille et mille applications, dont j’avais developpé la theorie, dans la section 3\textsuperscript{me} du II.\textsuperscript{livre} de cet ouvrage[$\dagger$)], du moins en allemand, beaucoup avant d’avoir vu l’ouvrage de Mr. \textsc{Legendre} — je dis, pouvais je parler de ce principe, que j’avais annoncé a plusieurs de mes amis deja en 1803 comme devant faire partie de l’ouvrage que je preparois, — comme d’une \textls[120]{methode empruntée} de Mr. \textsc{Legendre}? Je n’avait pas l’idée, que
\par\medskip\noindent\rule{.24\linewidth}{.4pt}\par\smallskip
{\small
\noindent\hspace*{\parindent}[\textsuperscript{*}) Es handelt sich hier um die Aufzeichnung des \textit{Tagebuchs} Nr. 88 vom 17. Juni 1798; siehe auch den Brief an \textsc{Olbers} vom 24. Januar 1812, Werke VIII, S. 140. Vergl. zum folgenden überhaupt die Werke VIII, S. 136—141 abgedruckten Veröffentlichungen und Briefstellen, wo übrigens S. 138 und S. 139 das Jahr 1794 — nicht wie im Text und in der \textit{Theoria motus}, art. 186, Werke VII, S. 253, 1795 — als dasjenige angegeben ist, seitdem \textsc{Gauss} die Methode der kleinsten Quadrate angewandt hatte.]
\par\noindent\hspace*{\parindent}[\textsuperscript{**}) \textit{Theoria motus corporum coelestium}, Hamburgi 1809, Werke VII (1906), S. 1.]
\par\noindent\hspace*{\parindent}[\textsuperscript{***}) \textsc{A. M. Legendre}, \textit{Nouvelles Méthodes pour la détermination des orbites des comètes}, Paris 1806, datiert vom 6. März 1805; \textsc{Abbe} hat 1871 noch auf eine zweite vor der \textit{Theoria motus} erschienene Arbeit verwiesen, nämlich auf die des Amerikaners \textsc{Adrain}: \textit{Research concerning the probabilities of the errors wich happen in making observations}, The Analyst or Mathem. Museum \textbf{1}, 1808, S. 93, der unabhängig von \textsc{Legendre} zur Methode der kleinsten Quadrate gelangt sein soll.]
\par\noindent\hspace*{\parindent}[$\dagger$) A. a. O. artt. 175—186, Werke VII (1906), S. 240—254.]
\par}
\clearpage\SourceLeaf{381}\ArithmeticHead{374}{L}{ANALYSIS.\quad BRIEFWECHSEL.}
\noindent Mr. \textsc{Legendre} pouvait attacher tant de prix à une idée aussi \textls[120]{simple}, qu’on doit plutot s’etonner qu’on ne l’a pas eue depuis 100 ans, pour se facher que je raconte, que je m’en suis servi avant lui? En effet il serait tres facile de le prouver à tout le monde par des temoignages qu’on ne saurait refuser, si cela valait la peine. Mais j’ai cru que tous ceux qui me connaissent le croiroient meme sur ma parole, ainsi que je l’aurait cru de tout mon coeur si Mr. \textsc{Legendre} avait avancé, qu’il avait possedé la methode deja avant 1795. J’ai dans mes papiers beaucoup de choses, donc peut être je pourrai perdre la priorité de la publication: mais soit, j’aime mieux faire \textls[120]{mûrir} les choses.
\par\noindent\makebox[\linewidth][s]{{.}\hfill{.}\hfill{.}\hfill{.}\hfill{.}\hfill{.}\hfill{.}\hfill{.}\hfill{.}\hfill{.}\hfill{.}\hfill{.}\hfill{.}\hfill{.}\hfill{.}\hfill{.}\hfill{.}\hfill{.}\hfill{.}\hfill{.}\hfill{.}\hfill{.}\hfill{.}\hfill{.}\hfill{.}\hfill{.}\hfill{.}}\par
\indent Continués, Monsieur, de m’honorer de votre bienveillance, que je range parmi les choses les plus essentielles à mon bonheur.
\par\noindent\hfill\textsc{Ch. Fr. Gauss.}\par
